% !TeX program = xelatex
\documentclass[a4paper,12pt]{ctexart}

% 导入自定义样式包
\usepackage{../../style/guilin-paper}
\usepackage{titletoc} 
\usepackage[hidelinks]{hyperref} % <--- 修复报错：添加此包以支持 \pdfbookmark

%===========================================================
%   Chinese TSD 标准目录样式定义
%===========================================================
\titlecontents{section}[0em]{\addvspace{5pt}\heiti\zihao{-4}}
    {\thecontentslabel\enspace}{}
    {\titlerule*[5pt]{.}\contentspage}

\titlecontents{subsection}[2em]{\addvspace{2pt}\songti\zihao{5}}
    {\thecontentslabel\enspace}{}
    {\titlerule*[5pt]{.}\contentspage}

%===========================================================
%   信息填写区域
%===========================================================
\courseName{自然辩证法概论}
\teacherName{孙\quad 伟}
\paperTitleLineA{周期性课题守恒论视角下}
\paperTitleLineB{科技异化的本质与超越}
\studentName{伍祺峻}
\studentID{2120251293}
\studentCollege{计算机科学与工程学院}
\studentGrade{2025 级}
\studentLocation{广西桂林 541006}
\studentPhone{13054083944} 
\studentEmail{1571379055@qq.com}
\studentBirthYear{2000}
\studentOrigin{湖南郴州}
\studentResearch{数据分析与应用}

%===========================================================
%   正文开始
%===========================================================
\begin{document}

% 1. 生成封面
\makeCoverPage

% 2. 论文首部信息
\begin{center}
    {\zihao{3}\heiti \getPaperTitleLineA \getPaperTitleLineB}
\end{center}

\begin{center}
    \vspace{-0.5em}
    \setlength{\baselineskip}{23pt} 
    {\zihao{-4}\kaishu \getStudentName} \\
    {\zihao{-4}\kaishu （桂林理工大学 \getStudentCollege，\getStudentLocation）}
\end{center}

\vspace{1em}

% 3. 摘要与关键词
\begin{PaperAbstract}{
科技异化，正如很多中国式关系的纠葛一样，是现代文明进程中一种既普遍又复杂的阵痛。它深刻折射出人与自然、人与自我之间那层难以亲近却又渴望和谐的张力。面对这一困境，单纯的技术批判往往流于表面。本文坚持马克思科学主义社会观，在历史唯物主义的框架下，深度融合斯多葛学派的客观宿命感与阿德勒心理学的能动目的论，创新性地构建“周期性课题守恒论”。我主张，“异化”并非文明走错了路，而是人类在背负着历史行囊前行时，确立主体性所必须经历的、带有某种必然色彩的“必修课题”。根据“体验守恒”机制，任何试图逃避的技术债，终将以某种其他形式进行重现。其中，解决之道不在于逃避，而在于我们如何去面对、感悟这一课题、去直面课题，寻找出最有用的核心原理，实现从必然王国向自由王国的跨越。
}
马克思科学主义社会观；周期性课题守恒；科技异化；斯多葛学派；阿德勒心理学
\end{PaperAbstract}

%===========================================================
%   4. 标准目录 (Chinese TSD Style)
%===========================================================
\newpage
\pagenumbering{Roman} % 目录页码使用罗马数字
\pdfbookmark[1]{目录}{anchor} % 现在不会报错了
\begin{center}
    \zihao{3}\heiti 目\quad 录
\end{center}
\vspace{1em}
\makeatletter
\@starttoc{toc}
\makeatother
\newpage
\pagenumbering{arabic} % 正文恢复阿拉伯数字，从1开始计数

%===========================================================
%   5. 正文内容
%===========================================================
\zihao{5}\songti
\setlength{\baselineskip}{23pt}
\setlength{\parindent}{2em} 

\section{引言：时代浪潮下的阵痛与纠葛}

首先，就如很多看似对立的关系一样，人类与科技之间的纠葛，也是一种既普通又复杂的纠结。为何说普通？因为自工业革命至今，我们所见所闻的历史大多如此——我们不断发明工具来解放自我，却又不断陷入“工具将要取代主人”的恐慌之中。


然而，在工具理性不断膨胀的今天，科技这把“双刃剑”正展现出它最残酷的一面。正如靳敏所言，它在递给我们物质福祉的同时，也无声无息地剥离了人的主体性与自然的本真状态\upcite{1}。宏观上，人工智能的巨轮正碾过旧有的社会结构；微观上，每一个个体的职业路径都充满了某种不可避免的焦虑。马克思认为，科技异化并非技术本身的“坏”，而是特定社会关系的底色所致\upcite{6}。基于此，我尝试引入斯多葛的客观智慧与阿德勒的能动逻辑，构建“周期性课题守恒论”。在我看来，这种关于“被替代”的焦虑，是文明进阶时必须经历的阵痛，我们不能像鸵鸟一样把头埋进沙子里，而必须在马克思主义的指引下，通过深刻的“体验”去完成这门历史的结业课。

\section{理论基石：马克思科学主义社会观的指引}

“周期性课题守恒论”并非无源之水，它深深植根于马克思关于“自然的人化”与“人的自然化”这一宏大叙事中。这为我们理解科技那层复杂的“外衣”提供了根本的本体论依据。

\subsection{异化的社会历史根源}
马克思曾一针见血地指出，"异化劳动"是所有问题的源头。在资本逻辑的剧本里，科技被异化成了掠夺的工具，不再是人类温情的延伸\upcite{6}。这表明，异化并非永恒的宿命，而是人类历史在特定的阶段，为了走向更成熟而必须付出的代价。

在《1844年经济学哲学手稿》中，马克思系统阐述了异化劳动的四个维度：人同自己的劳动产品相异化、人同自己的劳动活动相异化、人同自己的类本质相异化、人同人相异化。这四重异化在当代科技语境下呈现出新的表现形式。程序员编写的代码成为资本增殖的工具，而非自我表达的载体；设计师的创意被算法标准化，失去了独特性；知识工作者的劳动成果被平台垄断，个人价值被数据量化。

更深层次地看，杨英姿指出，人与自然关系的演进遵循唯物史观的内在逻辑\upcite{2}。从原始社会的"天人合一"，到农业文明的"人定胜天"，再到工业文明的"征服自然"，每一次关系转变都伴随着生产力的飞跃和新的矛盾积累。当前的生态危机和技术焦虑，正是这种矛盾在新时代的集中爆发。我们不能天真、简单地认为回到前工业时代就能解决问题，唯一的出路是在更高层次上实现人与自然、人与技术的辩证统一。

\subsection{科技的双重底色与逻辑内核}
在马克思的视野里，科技既是革命性的生产力，也承载着沉重的社会属性。承认这种矛盾，不仅是方法论的要求，更是一种看待世界的正确心态。这种“两点一线”的辩证逻辑，正是构建“周期性课题守恒论”的基石。

一方面，科技是人类解放的强大力量。它延伸了人的感官，放大了人的能力，缩短了必要劳动时间，为人的全面发展创造了物质前提。正如马克思所预见的，机器大工业的发展最终将为共产主义社会奠定基础。人工智能、生物技术、量子计算等前沿科技，都在以前所未有的速度推动生产力的跃升。

另一方面，在资本主义生产关系中，科技又成为加剧剥削和控制的手段。算法监控让劳动者的每一分钟都被精确量化，平台经济将灵活用工推向极致，大数据技术使得个人隐私无处遁形。技术进步非但没有缩短工作时间，反而通过"996"、"007"等形式强化了劳动强度。这种悖论的根源不在技术本身，而在于技术的所有权和使用权被少数人垄断。

\section{理论建构：周期性课题守恒论的逻辑展开}

在马克思主义的宏大视角下，我也试图融合心理学知识，去思考、拆解文明演进中那些微观的、若隐若现的机制。

\subsection{斯多葛视角：接纳历史的“逻各斯”}
斯多葛学派相信宇宙万物运行中存在着一种客观的“逻各斯”（Logos）\upcite{4}。对应到我的理论中，这一阶段的异化痛苦，正如人在锻炼中那些无法躲避的肌肉酸痛、疲惫一样，是必然的。顺应这种规律，意味着我们要先学会接纳。任何试图绕过这一阵痛阶段直接拥抱自由的念头，都显得过于幼稚且不可行。

然而，接纳并不意味着消极等待。斯多葛哲学强调的是一种"积极的顺应"：在认清客观规律的基础上，调整自己的心态和行动策略。当我们明白技术冲击是文明演进的必然环节，焦虑就会转化为清醒的准备。这种心态上的转变，是走出异化困境的第一步。正如爱比克泰德所言："不是事情本身困扰我们，而是我们对事情的看法困扰我们。"

从历史长河来看，每一次重大技术变革都伴随着阵痛期。工业革命初期的卢德运动，本质上是手工业者对机器的恐惧；电气化时代，人们担心电力会导致大规模失业；计算机普及之初，打字员、会计等职业面临威胁。但历史证明，技术进步最终创造了更多新的就业机会和生活方式。当前对人工智能的焦虑，或许在几十年后回望，也不过是又一次必经的阵痛。

\subsection{阿德勒视角：能动的“课题”与“体验守恒”}
如果说斯多葛告诉我们要接纳，那么阿德勒心理学则赋予了我们改变的勇气。阿德勒主张“目的论”，拒绝深陷于“原因论”的泥潭\upcite{5}。我基于此提出“体验守恒定律”：一个文明阶段所需的实践总量是守恒的。如果我们因为贪图安逸而选择在技术伦理探讨上偷懒，这些未被消化的体验迟早会化作历史的债务。

"体验守恒定律"的逻辑是：文明的每一次跃升，都需要集体意识完成相应的认知升级。这个过程无法跳过，只能延后。如果我们在工业革命时期逃避了对机器伦理的思考，那么这些问题会在信息革命时期加倍显现。如果我们在互联网时代回避了对算法权力的讨论，那么在人工智能时代就会面临更严重的失控风险。

具体而言，"体验守恒"包含三个层面：

其一，认知层面的守恒。每一代人都需要亲身经历技术冲击，形成自己的理解框架。父辈的经验无法直接传递，因为技术形态已经变化。这就是为什么每一代人都要重新讨论"技术是福是祸"这个古老命题。

其二，情感层面的守恒。技术焦虑不会因为被压抑而消失，它只会以其他形式爆发——可能是职场内卷，可能是教育焦虑，可能是社会撕裂。只有正视并处理这些情感，才能真正走出阴影。

其三，实践层面的守恒。技术素养的提升需要大量的实践积累。如果一个社会选择让少数精英掌握技术，而大众保持无知，那么最终会付出民主赤字、数字鸿沟扩大的代价。每个人都需要完成属于自己的"技术启蒙"。


\section{现象剖析：逃避课题导致的双重异化危机}

在“周期性课题守恒论”看来，当前社会的种种危机，本质上是因为我们在资本的诱惑下，长期试图走捷径，导致了严重的“科技负债”。

\subsection{宏观阵痛：技术变革下的集体焦虑}
如今科技变革带来的阵痛，正是那些被我们逃避的课题在集中爆发。以生成式人工智能（AIGC）为例，那种“被替代”的群体性恐慌，其实是我们过去长期过度依赖工具理性、习惯于将思考直接丢给算法所招致的惩罚。

ChatGPT的横空出世，让无数内容创作者、翻译工作者、客服人员陷入恐慌。但冷静分析会发现，真正被威胁的是那些机械性、重复性的劳动。那些仅仅依靠模板化写作、套路化翻译的从业者，本来就处于被替代的边缘。人工智能只是加速了这一进程，让问题提前暴露。

从更深层次看，这种集体焦虑反映了教育体系的滞后。我们的教育长期以来强调标准答案和应试技巧，而恰恰这些能力最容易被人工智能超越。如果教育体系能够更早地转向批判性思维、创造性解决问题、跨学科整合等"人之为人"的核心能力培养，今天的焦虑就会小得多。

\subsection{微观困境：职业发展中的主体性迷失}
作为一名计算机专业的学生，我对这种异化感的体会更为具体。在职业发展中，技术栈的快速迭代让我们焦虑。武春芳在解读《心灵捕手》时提到，人的心理防御往往源于缺乏直面过去的勇气\upcite{3}。如果我们只把代码看作糊口的工具，而不是自我实现的延伸，那么每次技术更新都会导致一次新的令人恐惧的事物诞生。

\section{路径重构：走向辩证复归的“结业”}

解决异化的途径，绝不是极端的砸碎机器这般简单，而是学会像成熟的成年人一样，背负着行囊穿透异化，去寻找那个名为“自由王国”的终点。


\subsection{变"自在"为"自为"的实践转化}

我们需要在职业和生命中重塑"目的论"。作为未来的技术从业者，我意识到，要打破这层科技带来的恐惧，我们必须将"被动适应"转变为"主动驾驭"。在害怕被替代的恐惧中寻找我们作为一个活生生的人的创造性的内核，通过主动的感受和体验去偿还那些我们因为没有思考带来的"科技欠债"。

马克思区分了"自在的阶级"和"自为的阶级"——前者只是客观存在的群体，后者则具有了自觉的阶级意识和行动能力。在技术语境下，我们同样需要从"自在的技术使用者"转变为"自为的技术主体"。

\subsection{从“自在”到“自为”的具体转变层面}
具体而言，这种转变包含以下几个层面：

首先，从技能学习到思维培养。不要满足于会用某个工具、掌握某个框架，而要理解背后的原理和思想。学习编程语言，更重要的是培养计算思维；使用人工智能，更重要的是理解算法逻辑和数据伦理。当我们拥有了这种"元能力"，技术更新就不再是威胁，而是扩展能力边界的机会。

其次，从被动接受到主动选择。面对层出不穷的新技术，我们要有选择性地投入精力。不是每个热点都值得追，不是每个趋势都要跟随。建立自己的判断标准，明确自己的兴趣方向，有意识地构建个人技术栈。这种主动性本身就是对异化的抵抗。

再次，从工具理性到价值理性。在技术实践中融入人文关怀和伦理考量。譬如说，作为程序员，不仅要思考"如何实现"，更要思考"是否应该实现"、"会带来什么影响"。技术不是价值中立的，每一行代码都包含着设计者的价值判断。当我们有意识地将伦理考量纳入开发流程，就是在重新夺回被异化的主体性。

最后，从个体奋斗到集体行动。技术异化不是个人努力就能完全解决的，它需要制度变革和社会运动的支持。参与开源社区、推动技术民主化、倡导算法透明，这些集体行动是对抗技术垄断和资本控制的有效方式。当技术从业者形成"自为的群体"，就有可能改变技术发展的方向。


\section{结语：完成这场漫长的和解}

综上所述，本文所构建的“周期性课题守恒论”，其核心逻辑在于：科技异化并非文明的偶然歧路，而是人类在确立主体性过程中，必须正面承接的一场“周期性课题”。

通过融合马克思主义的历史唯物主义、斯多葛学派的客观接纳和阿德勒心理学的能动目的论，本文试图提供一个既现实又富有希望的解决框架。我们需要斯多葛式的理性，去接纳那些客观存在的历史阵痛，明白异化是文明成长的必经行囊；我们更需要阿德勒式的勇气，将技术从单纯的"增殖工具"复归为"自我实现"的延伸。通过主动的变革与深刻的实践体验，去偿还科技债务，进而完成这一历史阶段的宿命课题。

这场和解的完成，需要几代人的共同努力。而我们这一代人的使命，就是认真地完成历史交给我们的这门课程，不逃避，不放弃，在实践中探索，在探索中成长，为后来者铺平道路。

\newpage
%===========================================================
%   6. 参考文献
%===========================================================
\begin{PaperBib}
    \item {[N]} 靳敏. 科学技术赋能人与自然和谐共生的路径探索. 江苏经济报, 2025-08-01(T07).
    \item {[J]} 杨英姿. 人与自然关系演进的唯物史观阐释. 哈尔滨工业大学学报(社会科学版), 2025, 27(5).
    \item {[J]} 武春芳. 人生如逆旅, 我亦是行人——基于阿德勒疗法对电影《心灵捕手》的心理学解读. 名作欣赏, 2024(21).
    \item {[M]} (古罗马)马可·奥勒留. 沉思录. 何怀宏, 译. 北京: 中央编译出版社, 2008.
    \item {[M]} (奥)阿尔弗雷德·阿德勒. 自卑与超越. 曹晚红, 译. 北京: 作家出版社, 2016.
    \item {[M]} 马克思, 恩格斯. 马克思恩格斯选集(第3卷). 北京: 人民出版社, 2012.
\end{PaperBib}

% 7. 作者简介
\makeBio

\end{document}